\documentclass[12pt]{article}

\usepackage[brazil]{babel}
\usepackage[utf8]{inputenc}
\usepackage[T1]{fontenc}
\usepackage{amsmath}
\usepackage{geometry}
\geometry{a4paper, margin=2.5cm}

\title{Avaliação de Matemática\Função do Primeiro Grau}
\author{Professor(a): \rule{8cm}{0.4pt}}
\date{Data: \rule{3cm}{0.4pt}}

\begin{document}

\maketitle

\vspace{0.5cm}

\textbf{Aluno(a):} \rule{10cm}{0.4pt} \
\textbf{Série/Ano:} \rule{4cm}{0.4pt}

\vspace{0.8cm}

\section*{Questão 1 – Conceito (1,0 ponto)}
Uma função do primeiro grau pode ser escrita na forma:
[
f(x) = ax + b
]
Explique com suas palavras o significado dos coeficientes (a) e (b).

\vspace{2cm}

\section*{Questão 2 – Identificação dos coeficientes (1,5 pontos)}
Dada a função:
[
f(x) = 3x - 5
]
Responda:

\begin{itemize}
\item a) Qual é o valor de (a)?
\item b) Qual é o valor de (b)?
\end{itemize}

\vspace{1.5cm}

\section*{Questão 3 – Cálculo de valores da função (2,0 pontos)}
Considere a função:
[
f(x) = 2x + 1
]
Calcule:

\begin{itemize}
\item a) (f(0))
\item b) (f(2))
\item c) (f(-1))
\end{itemize}

\vspace{1.5cm}

\section*{Questão 4 – Raiz da função (2,0 pontos)}
Determine a raiz da função:
[
f(x) = 4x - 8
]

\vspace{2cm}

\section*{Questão 5 – Situação-problema (3,5 pontos)}
Uma corrida de aplicativo cobra uma taxa fixa de R$ 5,00 mais R$ 2,00 por quilômetro rodado.

\begin{itemize}
\item a) Escreva a função que representa o valor da corrida em função da distância percorrida.
\item b) Quanto custará uma corrida de 10 km?
\end{itemize}

\vspace{2cm}

\begin{center}
\textbf{Boa prova!}
\end{center}

\end{document}

